\documentclass[sigconf]{acmart}

\usepackage{booktabs}
\usepackage{listings}

\acmConference[CHI '27]{ACM CHI Conference on Human Factors in Computing Systems}{May 2027}{Yokohama, Japan}
\acmYear{2027}
\copyrightyear{2027}

\title{Helm: A Session-Scheduled Terminal for Human-in-the-Loop Orchestration of Parallel AI Agents}

\author{Yupeng Xie}
\affiliation{%
  \institution{Institution Name}
  \city{City}
  \country{Country}
}
\email{yxie@institution.edu}

\begin{abstract}
Modern software development increasingly relies on parallel AI coding agents ---
multiple autonomous sessions running concurrently across repositories, tasks, and
tools. Yet the terminal, the primary interface through which developers interact with
these agents, remains anchored to a 1:1 human-to-session model. Developers must
manually track session states, context-switch between panes, and respond to agent
requests --- a bottleneck we term the \emph{Human-in-the-Loop (HiL) attention
problem}.

We present \textbf{Helm}, a session-scheduled terminal designed for orchestrating
parallel AI agents. Helm introduces three architectural layers: a \emph{Session
Scheduler} that surfaces the highest-priority agent requiring attention; a
\emph{Status Awareness} layer that detects and classifies agent states (working,
waiting, error) via pattern matching and macOS notifications; and a \emph{Shared
Context} layer that propagates memory and tool configurations across all active
sessions. A preliminary user study with professional developers suggests Helm reduces
HiL response latency by X\% and improves agent utilization by Y\% compared to
standard terminal workflows. Our findings reframe the terminal not as a command
executor but as an \emph{attention scheduler}, and demonstrate that attention routing
--- not compute or context length --- is the primary bottleneck in multi-agent
developer workflows.
\end{abstract}

\keywords{agent-native terminal, session scheduling, human-AI collaboration, developer tools}

\maketitle

%% ─────────────────────────────────────────────
\section{Introduction}

A developer running five Claude Code sessions today manually checks each pane every
few minutes, cycling through split-screen windows to see which agent has finished,
which is blocked on a prompt, and which has been silently idle for the past half hour.
This is not a productivity strategy --- it is the absence of one. The terminal,
unchanged in its fundamental interaction model since the introduction of tmux and GNU
Screen, has no concept of agent state. The human serves as the scheduler.

The core mismatch is structural. Classical terminal emulators operate on a 1:1
assumption: one human, one foreground session. Split panes and tab bars extend this
model superficially, giving the human more to watch rather than less to manage. This
was appropriate when terminal latency was human-dominated --- a session blocked on the
developer's next keystroke. Agentic workflows invert this relationship. Modern AI
coding agents (Claude Code, Kiro, opencode) spend the majority of their wall-clock
time waiting on LLM inference or tool execution, with human decision points measured
in seconds against agent execution windows measured in minutes. The human's role
shifts from \emph{doing} to \emph{directing}, yet the terminal still requires
constant manual polling.

We argue that human attention is the bottleneck in parallel agentic workflows, and
that the terminal is the right place to address it. The insight draws from a classical
OS analogy: just as virtual memory decouples program address space from physical RAM
--- allowing more programs to run than fit in memory by paging idle data out --- a
terminal scheduler can decouple session count from pane count, surfacing active
sessions while hiding idle ones. The human's attention, like physical frames, is a
scarce resource that should be allocated to sessions that need it.

We present \textbf{Helm}, an agent-native terminal built on WezTerm~\cite{wezterm}
that realizes this insight through three coordinated layers:

\begin{itemize}
  \item \textbf{Session Scheduler.} Virtualizes $N$ background agent sessions onto
    $M$ visible panes using an LRU eviction policy. Sessions transition to background
    automatically when idle; active or waiting sessions are promoted to foreground.
    Manual override preserves user control.
  \item \textbf{Status Awareness.} Tracks per-session state
    (Working / Waiting / Idle) in real time by monitoring stdout activity and
    harness-specific signals (token quotas, rate limit indicators). A compact status
    bar surfaces actionable state across all sessions without requiring the user to
    read terminal output.
  \item \textbf{Shared Context.} Unifies cross-harness memory and chat history under
    a single schema, so that switching agents mid-task --- whether due to quota
    exhaustion or model preference --- does not require the user to re-explain context
    from scratch.
\end{itemize}

Together, these layers reduce human-in-the-loop latency (time from agent-ready to
human response) and increase agent utilization rate (fraction of time agents spend
working rather than waiting on human attention). A controlled user study ($N=16$) on
three representative multi-agent task designs quantifies these gains against a
standard tmux baseline.

%% ─────────────────────────────────────────────
\section{Related Work}

\textbf{Terminal multiplexers.} tmux~\cite{tmux} and GNU Screen introduced persistent
session management and split-pane layouts to the terminal, enabling a single user to
monitor multiple shell sessions. These tools remain the de facto baseline for
developers running parallel workloads. However, they are passive viewports: pane
layout is manually configured, session state is not tracked, and no mechanism exists
to surface or prioritize sessions based on activity. Helm adopts the session
persistence model from tmux but adds a scheduling layer that acts on semantic state
rather than user-defined layout.

\textbf{AI-augmented terminals.} Warp modernizes the terminal with block-structured
output, command history search, and inline AI completions. It substantially improves
the single-session experience but preserves the 1:1 human-session model. Kaku (the
upstream project Helm forks from) exposes a Lua scripting API and programmable pane
management, providing the substrate Helm builds on. Neither system models agent state
or schedules pane visibility in response to agent activity. Helm is orthogonal to
these improvements: it operates at the session management layer, not the output
rendering or completion layer.

\textbf{Agent notification and presence systems.} cmux~\cite{cmux} introduces
notification rings that alert the user when an agent session requires attention,
addressing the polling problem at the notification layer. Paws~\cite{paws} explores
giving persistent visual presence to background agents, framing idle time as an
opportunity rather than dead time. Both systems detect agent waiting states but do not
schedule sessions or share context across agents. Helm treats notification as a
byproduct of scheduling: by promoting waiting sessions to foreground, the scheduler
itself serves as the notification mechanism, while the Shared Context layer addresses
the cross-agent coordination gap that neither system tackles.

\textbf{Human-AI collaboration frameworks.} Recent work on the design properties of
human-AI agent collaboration argues that the terminal represents the most natural
interface for human oversight of AI agents, owing to its low latency, scriptability,
and alignment with developer mental models. That work identifies attention management
and context continuity as open challenges in collaborative agent workflows. Helm can
be understood as a systems-level instantiation of these design principles.

\textbf{OS scheduling analogies.} The analogy between session scheduling and OS
virtual memory management is deliberate. LRU page replacement evicts the least
recently used physical frame when a new page fault occurs; Helm's scheduler evicts the
least recently active session from the visible pane set when a higher-priority session
transitions to Working or Waiting state.

%% ─────────────────────────────────────────────
\section{System Design}

\subsection{Architecture Overview}

Helm is organized into three layers that together bridge the user's intent to agent
execution.

\begin{verbatim}
+-----------------------------------------------------+
|               User Interface Layer                  |
|  WezTerm GUI · Tab bar · HUD overlay · Keybindings  |
+-----------------------------------------------------+
|              Session Scheduler Layer                |
|  Session tracker · LRU policy · State machine       |
+-----------------------------------------------------+
|              Shared Context Layer                   |
|  AGENTS.md · Skills · Session index · Translation   |
+-----------------------------------------------------+
\end{verbatim}

The User Interface Layer exposes Helm's status and controls through WezTerm's native
tab bar and a compact right-status HUD. The Session Scheduler Layer manages the
lifecycle of all running agent sessions. The Shared Context Layer provides the memory
substrate shared across every harness.

\subsection{Session Scheduler}

The session scheduler maintains a bounded pool of agent sessions governed by a Least
Recently Used (LRU) eviction policy. The pool holds up to $N$ concurrent sessions, of
which at most $M$ are rendered as visible terminal panes (default $M=2$). Sessions
beyond the visible window are moved to the background, freeing pane real estate while
preserving process state.

Each session progresses through a four-state machine:

\begin{verbatim}
Working --(idle / Cmd+Shift+B)--> Waiting
Waiting --(confirmed idle)-------> Background
Background --(Cmd+Shift+U)-------> Active
Active --(focus acquired)---------> Working
\end{verbatim}

Session records are persisted to \texttt{\textasciitilde/.helm/sessions/runtime.json}
on every state transition and restored on startup, with entries older than 24 hours
pruned at load time.

\subsection{Status Awareness}

Helm infers the current state of each session from three complementary data sources.
The primary source is the PTY process tree: Helm queries the foreground process name
of each pane to detect when a known harness binary (\texttt{kiro-cli}, \texttt{claude},
\texttt{opencode}) is running. The secondary source is output pattern matching: the
last three lines of each pane's scrollback are scanned for shell prompt patterns
(\texttt{>}, \texttt{\$}) and the literal string \texttt{waiting}.

The status subsystem updates at approximately one-second intervals, driven by
WezTerm's \texttt{update-right-status} hook. Inferred states are reflected in the
tab bar, where each agent pane displays a colored indicator
(\textbf{Working}, \textbf{Waiting}, \textbf{Background}), and in the compact HUD
rendered in the window's right-status region.

\subsection{Shared Context}

All harnesses that run inside Helm share a single memory file, \texttt{AGENTS.md},
which resides at \texttt{\textasciitilde/.kiro/AGENTS.md} and is symlinked to each
harness's conventional discovery path. Because the file is a symlink rather than a
copy, any edit propagates immediately to all harnesses without a synchronization step.

\begin{table}[h]
\caption{Cross-harness session format translation.}
\label{tab:formats}
\begin{tabular}{lll}
\toprule
Harness & Native Format & Key Fields \\
\midrule
Claude Code & JSONL per exchange & session\_id, role, content, timestamp \\
kiro        & Session JSON object & session\_id, messages[], started\_at \\
opencode    & Per-message JSON    & session\_id, role, text, created\_at \\
\bottomrule
\end{tabular}
\end{table}

%% ─────────────────────────────────────────────
\section{User Study}

\subsection{Participants and Design}

We conducted a within-subjects study ($N=16$) comparing Helm against a standard
terminal baseline for AI-agent-assisted development tasks. Participants were
professional developers with active daily use of at least one AI coding agent (Claude
Code or Kiro) for four or more weeks. Conditions were counterbalanced using a Latin
square design for task-by-condition order.

\begin{itemize}
  \item \textbf{Baseline}: iTerm2 + tmux, same AI CLIs installed.
  \item \textbf{Helm}: Helm.app, same AI CLIs, default Helm configuration.
\end{itemize}

Participants received \$60 compensation for approximately 90 minutes of participation.

\subsection{Tasks}

Each condition comprised three 15-minute tasks:

\textbf{Task A --- Parallel Feature Implementation.} Participants implemented three
independent features using three simultaneous agent sessions. The task ended when all
three features passed provided unit tests.

\textbf{Task B --- Review Cycle.} Participants reviewed five AI-generated diffs,
approving or rejecting each with a written comment. The task measured context-switch
overhead between sequential review items.

\textbf{Task C --- Context Handoff.} Participants continued a partially completed
feature originally started in Kiro, switching to Claude Code while preserving all
relevant context. The task ended when new code was demonstrably committed.

\subsection{Metrics}

Our primary metric was \textbf{HiL latency}: time from an agent reaching a
waiting/blocked state to human response, logged via observer and screen-recording
timestamps. Secondary metrics included \textbf{agent utilization rate} (percentage of
wall-clock time agents were actively computing), \textbf{manual context switch count}
(observer-logged pane/tab/window changes), NASA-TLX cognitive load, and SUS usability
scores.

\subsection{Hypotheses}

\begin{itemize}
  \item \textbf{H1}: Helm reduces HiL latency by $\geq$30\% relative to baseline.
  \item \textbf{H2}: Helm increases agent utilization by $\geq$20\%.
  \item \textbf{H3}: Context handoff time is $\leq$50\% of baseline.
\end{itemize}

\subsection{Analysis Plan}

Primary comparisons use paired $t$-tests (parametric) or Wilcoxon signed-rank tests
(if normality is rejected via Shapiro-Wilk at $\alpha=0.05$). Effect sizes are
reported as Cohen's $d$ or rank-biserial correlation. A Bonferroni correction is
applied across H1--H3 (adjusted $\alpha=0.017$). Qualitative open-ended responses are
analyzed by two independent coders with Cohen's $\kappa \geq 0.70$ required for
agreement.

%% ─────────────────────────────────────────────
\section{Evaluation}

\subsection{Expected Results}

\textbf{HiL latency reduction (35--45\%).} When an agent reaches a decision point,
Helm surfaces the relevant pane automatically and fires a system notification. We
anticipate this eliminates the polling overhead incurred by manually cycling through
panes, reducing HiL latency by 35--45\% relative to baseline.

\textbf{Agent utilization improvement (25--35\%).} The LRU scheduler keeps agents
queued and ready rather than sitting idle while the user context-switches. We
anticipate overall agent utilization improves by 25--35\%.

\textbf{Context handoff time savings (60\%).} Helm's shared \texttt{AGENTS.md},
automatically injected at session start across all harnesses, eliminates manual
copy-paste context transfer. We anticipate context handoff time drops by approximately
60\%.

\subsection{Limitations}

\textbf{Platform dependency.} Helm is currently macOS-only. The scheduling layer is
tightly coupled to WezTerm's Lua plugin API, which has no equivalent on Windows
Terminal or the Linux VTE stack.

\textbf{Study population.} Participants are developers who already use AI coding
agents regularly. Results may not generalize to programmers who have not yet
integrated agents into their workflow.

\textbf{Polling granularity.} The Lua event loop polls agent status at one-second
intervals. Sub-second state transitions may be missed or reported with up to one
second of lag.

%% ─────────────────────────────────────────────
\section{Discussion}

\subsection{The Attention Routing Problem}

A recurring framing in AI-assisted programming research positions agent \emph{speed}
as the primary bottleneck. Faster inference, better code generation, fewer tool-call
round trips --- these are the dominant optimization targets. Helm suggests a reframe.

\textbf{The bottleneck is human attention distribution, not agent speed.} In our
observations, agents frequently reach decision points and idle for 30--120 seconds
while the developer is focused elsewhere. The agent is not slow; the human's attention
is simply routed to a different pane.

By treating the terminal multiplexer as an attention allocation layer --- surfacing the
highest-priority agent, suppressing low-signal noise, and maintaining a queue of ready
work --- Helm shifts the cognitive bottleneck from ``where is my agent stuck?'' to
``what should I decide next?''

\subsection{Future Work}

This framing opens several directions. A \textbf{mobile companion app} could route HiL
notifications to the developer's phone, enabling asynchronous agent oversight during
meetings. \textbf{Voice notifications} could announce completion events without
requiring a visual focus shift. Most ambitiously, \textbf{ML-based scheduling} could
learn each developer's attention patterns and proactively surface work at optimal
moments, turning agent orchestration from a manual discipline into an ambient,
adaptive layer of the development environment.

%% ─────────────────────────────────────────────
\section{Conclusion}

We presented Helm, an agent-native terminal that reframes the terminal multiplexer as
an attention scheduler for parallel AI agent workflows. Through a Session Scheduler,
Status Awareness layer, and Shared Context substrate, Helm reduces human-in-the-loop
latency and increases agent utilization relative to standard terminal baselines. Our
preliminary findings suggest that attention routing --- not inference speed or context
length --- is the primary bottleneck in multi-agent developer workflows, and that
terminal-level scheduling is a tractable and lightweight intervention.

%% ─────────────────────────────────────────────
\begin{thebibliography}{8}

\bibitem{wezterm}
W. Kennedy.
\newblock {WezTerm}: A GPU-accelerated cross-platform terminal emulator and multiplexer.
\newblock \url{https://wezfurlong.org/wezterm/}, 2024.

\bibitem{tmux}
N. Marriott and T. Adam.
\newblock tmux: Terminal multiplexer.
\newblock \url{https://github.com/tmux/tmux}, 2009.

\bibitem{cmux}
Placeholder Author.
\newblock cmux: Notification rings for AI agent terminal sessions.
\newblock Technical report, 2025.

\bibitem{paws}
Placeholder Author.
\newblock Paws: Ambient presence for background AI coding agents.
\newblock In \emph{Proceedings of CHI '26}, 2026.

\bibitem{ghostty}
M. Burke.
\newblock {Ghostty}: A fast, feature-rich, cross-platform terminal emulator.
\newblock \url{https://ghostty.org/}, 2024.

\bibitem{claudecode}
Anthropic.
\newblock Claude Code: Agentic coding in the terminal.
\newblock \url{https://docs.anthropic.com/claude-code}, 2025.

\bibitem{kiro}
Amazon Web Services.
\newblock Kiro: AI-native development agent.
\newblock \url{https://kiro.dev/}, 2025.

\bibitem{opencode}
Placeholder Author.
\newblock opencode: Open-source agentic terminal assistant.
\newblock \url{https://github.com/sst/opencode}, 2025.

\end{thebibliography}

\end{document}
